\section*{第 \ref{ch:g5:series-parallel-circuits} 章——串联和并联电路}

\begin{solution}{exo:g5:series-parallel-circuits:1}
\omterm{def:g5:series-parallel-circuits:parallel}{支路}是电路分出的各条独立的
路径之一；\omterm{def:g5:series-parallel-circuits:parallel}{节点}是\omterm{def:g5:series-parallel-circuits:parallel}{支路}分开或汇合的地方。两个\omterm{def:g3:first-electric-circuit:bulb}{灯泡}
\omterm{def:g5:series-parallel-circuits:parallel}{并联}时，各自位于一条带\omterm{def:g3:first-electric-circuit:bulb}{灯泡}的\omterm{def:g5:series-parallel-circuits:parallel}{支路}上。
\end{solution}

\begin{solution}{exo:g5:series-parallel-circuits:2}
独立性（断开只影响自己的\omterm{def:g5:series-parallel-circuits:parallel}{支路}）；亮度充足
（每个\omterm{def:g3:first-electric-circuit:bulb}{灯泡}都和单独接入时一样亮，\omterm{def:g3:first-electric-circuit:battery}{电池}因此耗电更快）；控制点
（\omterm{def:g5:series-parallel-circuits:parallel}{支路}\omterm{ex:g3:first-electric-circuit:switch}{开关}只控制本\omterm{def:g5:series-parallel-circuits:parallel}{支路}，
分岔前的\omterm{ex:g3:first-electric-circuit:switch}{开关}控制全部）。
\end{solution}

\begin{solution}{exo:g5:series-parallel-circuits:3}
\omterm{def:g3:first-electric-circuit:bulb}{灯泡} 2 不受影响，仍然亮着——它自己的路径是通的。\omterm{def:g4:series-circuits:series}{串联}时
它却熄灭了，因为两个\omterm{def:g3:first-electric-circuit:bulb}{灯泡}共用同一条路，
任何一处断开都会同时使两个\omterm{def:g3:first-electric-circuit:bulb}{灯泡}断路。
\end{solution}

\begin{solution}{exo:g5:series-parallel-circuits:4}
B 组（\omterm{def:g5:series-parallel-circuits:parallel}{并联}）每个\omterm{def:g3:first-electric-circuit:bulb}{灯泡}都更亮——各自达到
满亮度。B 组\omterm{def:g3:first-electric-circuit:battery}{电池}耗电也更快：
它同时为两条\omterm{def:g5:series-parallel-circuits:parallel}{支路}供电。
\end{solution}

\begin{solution}{exo:g5:series-parallel-circuits:5}
两个\omterm{ex:g3:first-electric-circuit:switch}{开关}都闭合：两个\omterm{def:g3:first-electric-circuit:bulb}{灯泡}都亮。干路\omterm{ex:g3:first-electric-circuit:switch}{开关}断开：两个\omterm{def:g3:first-electric-circuit:bulb}{灯泡}都灭，
因为它控制的是公共路段。干路\omterm{ex:g3:first-electric-circuit:switch}{开关}闭合、\omterm{def:g5:series-parallel-circuits:parallel}{支路}\omterm{ex:g3:first-electric-circuit:switch}{开关}断开：\omterm{def:g3:first-electric-circuit:bulb}{灯泡} 1 亮，
\omterm{def:g3:first-electric-circuit:bulb}{灯泡} 2 灭。
\end{solution}

\begin{solution}{exo:g5:series-parallel-circuits:6}
若是\omterm{def:g4:series-circuits:series}{串联}，一个坏掉的\omterm{def:g3:first-electric-circuit:bulb}{灯泡}（或一盏被关掉的灯！）就会
使整座房子变暗，而且灯越多，所有灯就越暗。
\omterm{def:g5:series-parallel-circuits:parallel}{并联}接线带来的便利：每个房间的\omterm{ex:g3:first-electric-circuit:switch}{开关}只控制
本房间；厨房\omterm{def:g3:first-electric-circuit:bulb}{灯泡}烧坏，其余房间的灯仍然亮着；
每盏灯都能保持满亮度。
\end{solution}

\begin{solution}{exo:g5:series-parallel-circuits:7}
位于分岔之前——在电流进入住宅的公共路段上，
即保险丝盒中的总\omterm{ex:g3:first-electric-circuit:switch}{开关}。它能同时切断整座房子的电路，
用于检修和紧急情况。
\end{solution}

\begin{solution}{exo:g5:series-parallel-circuits:8}
\omterm{def:g5:series-parallel-circuits:parallel}{并联}时：关掉天花板灯不会熄灭床头灯，
拔掉床头灯的插头也不会使天花板灯变暗——
两者各自位于自己的\omterm{def:g5:series-parallel-circuits:parallel}{支路}上。
\end{solution}

\begin{solution}{exo:g5:series-parallel-circuits:9}
用我们普通的\omterm{ex:g3:first-electric-circuit:switch}{开关}，这些设计总会失败：与\omterm{def:g3:first-electric-circuit:bulb}{灯泡}\omterm{def:g4:series-circuits:series}{串联}的\omterm{ex:g3:first-electric-circuit:switch}{开关}只能从一端断开；
两个普通\omterm{ex:g3:first-electric-circuit:switch}{开关}\omterm{def:g4:series-circuits:series}{串联}时，要求\emph{两个都}闭合
（任一端都能关灯，却不能总是开灯）；两个\omterm{ex:g3:first-electric-circuit:switch}{开关}\omterm{def:g5:series-parallel-circuits:parallel}{并联}时，要使灯灭则要求
两个都断开。走廊的便利需要一种\emph{能在两条路之间做选择}的\omterm{ex:g3:first-electric-circuit:switch}{开关}，
而不只是断开一条路——这就是
双控\omterm{ex:g3:first-electric-circuit:switch}{开关}，
走廊中成对使用的正是它。
\end{solution}

\begin{solution}{exo:g5:series-parallel-circuits:10}
好处是：\omterm{def:g3:first-electric-circuit:bulb}{灯泡}位于\omterm{def:g5:series-parallel-circuits:parallel}{并联}\omterm{def:g5:series-parallel-circuits:parallel}{支路}上，一个坏掉的\omterm{def:g3:first-electric-circuit:bulb}{灯泡}
只会使它自己变暗。警告是：每少一个\omterm{def:g3:first-electric-circuit:bulb}{灯泡}就少一条\omterm{def:g5:series-parallel-circuits:parallel}{支路}，
\omterm{def:g3:first-electric-circuit:battery}{电池}（或小彩灯串自带的小电源）会把全部输出
都送到剩下的\omterm{def:g5:series-parallel-circuits:parallel}{支路}上。\omterm{def:g3:first-electric-circuit:bulb}{灯泡}太少时，留下的\omterm{def:g3:first-electric-circuit:bulb}{灯泡}各自承受的
电流超过设计值，发热更厉害，
因而很快就烧坏了。
\end{solution}

\begin{solution}{exo:g5:series-parallel-circuits:11}
应接在\omterm{def:g5:series-parallel-circuits:parallel}{支路}\omterm{ex:g3:first-electric-circuit:switch}{开关}自己的那段路上：位于\omterm{def:g5:series-parallel-circuits:parallel}{支路}\omterm{ex:g3:first-electric-circuit:switch}{开关}
与\omterm{def:g3:first-electric-circuit:bulb}{灯泡} 2 之间（即在该\omterm{def:g5:series-parallel-circuits:parallel}{支路}与其余部分汇合之前）。它
恰好在该\omterm{def:g5:series-parallel-circuits:parallel}{支路}工作时才亮，并且与\omterm{def:g3:first-electric-circuit:bulb}{灯泡} 2 \omterm{def:g4:series-circuits:series}{串联}，
两者位于同一条\omterm{def:g5:series-parallel-circuits:parallel}{支路}上。
\end{solution}
